\begin{enumerate}
	\item
	\begin{enumerate}
		\item L'étendue est la différence entre la plus grande et la plus petite valeur de la série :

\[14,30 - 12,86 = 1,44.\]

L'étendue de la série est $1,44~^{\circ}$C.

		\item La formule est \texttt{=MIN(B2:K2)}.
	\end{enumerate}

	\item

\[10,73 \times \left(1 + \frac{25}{100}\right) \approx 13,41 \]
	
Si la température avait augmenté de $25\,\%$, elle serait de $13,41^{\circ}$C. Or, la température est de $13,50^{\circ}$C soit supérieure. L'élève a raison, l'augmentation est bien supérieure à $25\,\%$.

	\item De 2024 à 2034, il s'écoule 10 années :

\[13,50 + 10\times 0,025 = 13,50 + 0,25 = 13,75.\]

La température moyenne estimée en 2034 est bien de $13,75~^{\circ}$C.

	\item $2044 - 2024 = 20$.
	
La variable \og Nombre d'années d'écart \fg{} contient la valeur 20.
	\item
	\begin{enumerate}
		\item $13,5 +  0,025x = 15 \iff 0,025x = 1,5 \iff x = \dfrac{1,5}{0,025}=60$.
		
La solution de l'équation est $x = 60$.
		\item La question posée revient à résoudre l'équation $13,5+0,025x = 15$. D'après la question précédente $x = 60$ où $x$ désigne le nombre d'années écoulées depuis 2024.

$2024 + 60 = 2084$.

La température moyenne annuelle à Paris serait estimée à $15~^{\circ}$C en 2084.
	\end{enumerate}
\end{enumerate}
